\section*{Abstract}
		This paper systematically analyzes the parameter dependency between SI units and T0-model natural units, revealing that fundamental parameters like $\alpha_{\text{EM}}$ and Yukawa couplings have dramatically different numerical values in different unit systems. Through detailed calculations, we demonstrate that direct transfer of parameter values between systems leads to errors spanning multiple orders of magnitude. The analysis extends beyond specific parameters to establish universal transformation rules and provides critical warnings against naive parameter transfer. This work establishes that the apparent inconsistencies in T0-model parameters are actually systematic unit-system dependencies that require careful transformation protocols for experimental verification.

	

	\section{Introduction}
	\label{sec:introduction}
	
	\subsection{The Parameter Transfer Problem}
	\label{subsec:parameter_problem}
	
	The T0 model, formulated in natural units where $\hbar = c = G = k_B = \alpha_{\text{EM}} = \beta_{\text{T}} = 1$, presents a fundamental challenge when compared with experimental data expressed in SI units. This paper demonstrates that the apparent inconsistencies between T0-model predictions and experimental observations are not physical contradictions but systematic unit-system dependencies.
	
	The core insight is that parameters such as $\alpha_{\text{EM}}$ represent fundamentally different quantities when expressed in different unit systems:
	
	$$\alphaEMSI \neq \alphaEMnat$$
	
	The parameter $\xipar$, by contrast, is a ratio of lengths and has the same value in both systems (Section~\ref{subsec:xi_transformation}).
	
	\subsection{Scope and Methodology}
	\label{subsec:scope}
	
	This analysis covers:
	\begin{itemize}
		\item Systematic calculation of parameter ratios between SI and T0-natural units
		\item Demonstration of transformation invariance for dimensionless ratios
		\item Extension to variable parameters like $\xipar$ and Yukawa couplings
		\item Universal warnings against direct parameter transfer
		\item Guidelines for correct experimental comparison protocols
	\end{itemize}
	
	\section{The $\xipar$ Parameter: Variable Across Mass Scales}
	\label{sec:xi_parameter}
	\section{The Universal $\xi$-Field Framework}
	
	The cornerstone of the T0-model is the universal geometric constant that serves as the fundamental parameter for all physical calculations.

		The universal geometric constant:
		\begin{equation}
			\xi = \frac{4}{3} \times 10^{-4} = 1.3333... \times 10^{-4}
		\end{equation}

	This dimensionless constant is used throughout T0 theory to connect quantum mechanical and gravitational phenomena. It establishes the characteristic strength of field interactions and provides the foundation for unified field descriptions.

		For the detailed derivation and physical justification of this parameter, see the document ``Parameter Derivation'' (available at: ).

	This geometric constant determines a characteristic energy scale for the $\xi$-field:
	
	\begin{equation}
		E_\xi = \frac{1}{\xi} = \frac{3}{4 \times 10^{-4}} = 7500 \text{ (natural units)}
	\end{equation}
	\subsection{Definition and Physical Meaning}
	\label{subsec:xi_definition}
	
	The parameter $\xipar$ is also the ratio of the Schwarzschild radius to the Planck length:
	
	\begin{equation}
		\xipar = \frac{r_0}{\lP} = \frac{2Gm}{c^2\lP}
		\label{eq:xi_definition}
	\end{equation}
	
\textbf{Crucial:} The parameter $\xipar$ scales with the mass of the object under consideration according to $\xipar(m) = 2Gm/(c^2\lP)$. The Higgs mass defines the fundamental reference scale $\xipar_0 = 1.33 \times 10^{-4}$, to which all other masses are normalized in the T0 model.
	
	\subsection{Connection to Higgs Physics}
	\label{subsec:xi_higgs_connection}
	
	The T0 model connects $\xipar$ with Higgs sector physics through the following relation; it is a consistency check, not an exact derivation (Doc.~385).
	
	\begin{equation}
		\xipar \approx \frac{\lambdah^2 v^2}{16\pichar^3 m_h^2} \approx 1.30 \times 10^{-4}
		\label{eq:xi_higgs_fundamental}
	\end{equation}
	
	where:
	\begin{itemize}
		\item $\lambdah \approx 0.13$ (Higgs self-coupling)
		\item $v \approx 246$ GeV (Higgs VEV)
		\item $m_h \approx 125$ GeV (Higgs mass)
	\end{itemize}
	
	The Higgs expression lies about $2.3\,\%$ below the setting $\xipar = 4/30000 = 1.33 \times 10^{-4}$ (Doc.~385). It supports the universal scale parameter from Standard Model physics, while the mass-dependent form $\xipar = 2Gm/(c^2\lP)$ applies to specific objects.
	
	\subsection{$\xipar$ Values in the SI System}
	\label{subsec:xi_si_values}
	
	Using SI constants:
	\begin{align}
		G &= 6.674 \times 10^{-11} \text{ m}^3/(\text{kg} \cdot \text{s}^2) \\
		\lP &= 1.616 \times 10^{-35} \text{ m}
	\end{align}
	
	We calculate $\xipar_{\text{SI}}$ for various objects:
	
	\begin{table}[htbp]
		\centering
		\adjustbox{max width=\textwidth}{%
\begin{tabular}{lcc}
			\toprule
			\textbf{Object} & \textbf{Mass} & \textbf{$\xipar_{\text{SI}}$} \\
			\midrule
			Electron & $9.109 \times 10^{-31}$ kg & $8.37 \times 10^{-23}$ \\
			Proton & $1.673 \times 10^{-27}$ kg & $1.54 \times 10^{-19}$ \\
			Human (70 kg) & $7.0 \times 10^{1}$ kg & $6.4 \times 10^{9}$ \\
			Earth & $5.972 \times 10^{24}$ kg & $5.5 \times 10^{32}$ \\
			Sun & $1.989 \times 10^{30}$ kg & $1.8 \times 10^{38}$ \\
			Planck mass & $2.176 \times 10^{-8}$ kg & $2.0$ \\
			\bottomrule
		\end{tabular}%
}
		\caption{$\xipar$ values for different objects in SI units}
		\label{tab:xi_si_values}
	\end{table}
	
	\textbf{The parameter $\xipar$ varies over about 60 orders of magnitude!}
	
	\subsection{$\xipar$ in T0-Natural Units}
	\label{subsec:xi_transformation}
	
	With $\xipar = r_0/\lP = 2Gm/(c^2\lP)$, $\xipar$ is a pure ratio of lengths. In natural units $G = 1/E_P^2$ and $\lP = 1/E_P$, so again
	\begin{equation}
		\xipar = \frac{2m}{E_P} = \frac{2m}{m_P}.
	\end{equation}
	The numerical value is therefore the same in SI and in T0-natural units; no conversion factor between the systems arises for this quantity. The values in Table~\ref{tab:xi_si_values} hold unchanged in both systems.
	
	\subsection{Invariance of Ratios}
	\label{subsec:xi_ratio_invariance}
	
	\textbf{Critical verification:} Since the values themselves are already the same in both systems, this holds all the more for the ratios between different objects:
	
	\begin{equation}
		\frac{\xipar_{\text{Sun}}}{\xipar_{\text{e}}} = \frac{1.83 \times 10^{38}}{8.37 \times 10^{-23}} = \frac{m_{\text{Sun}}}{m_e} = 2.18 \times 10^{60}
	\end{equation}
	
	\boxed{\text{Ratios are invariant under system transformation!}}
	
\section{The Fine-Structure Constant $\alpha_{\text{EM}}$}
\label{sec:alpha_em}

\subsection{The Mystification of 1/137}
\label{subsec:alpha_mystification}

The fine-structure constant $\alpha_{\text{EM}} \approx 1/137$ has been declared one of the greatest mysteries of physics by prominent physicists:

\begin{itemize}
	\item \textbf{Richard Feynman}: ``It is one of the greatest damn mysteries of physics: a magic number that comes to us with no understanding whatsoever.''
	\item \textbf{Wolfgang Pauli}: ``When I die, I will ask God two questions: Why relativity? And why 137? I believe he will have an answer for the first one.''
	\item \textbf{Max Born}: ``If $\alpha$ were larger, molecules could not exist, and there would be no life.''
\end{itemize}

\subsection{Electromagnetic Duality as the Key}
\label{subsec:electromagnetic_duality}

What all these statements overlook: The fine-structure constant possesses two mathematically equivalent representations that reveal its true nature:

\begin{align}
	\alpha_{\text{EM}} &= \frac{e^2}{4\pi\varepsilon_0\hbar c} \quad \text{(Standard form)} \label{eq:alpha_standard}\\
	\alpha_{\text{EM}} &= \frac{e^2 \mu_0 c}{4\pi \hbar} \quad \text{(Dual form)} \label{eq:alpha_dual}
\end{align}

This equivalence is based on the Maxwell relation $c^2 = \frac{1}{\varepsilon_0\mu_0}$ and reveals a fundamental electromagnetic duality:

\begin{equation}
	\frac{1}{\varepsilon_0 c} = \mu_0 c
	\label{eq:em_duality}
\end{equation}

\subsection{The Dual Nature of $\alpha$: System-Dependent yet Invariant}
\label{subsec:double_nature}

The fine-structure constant possesses a remarkable dual nature:

\subsubsection{As an Invariant Ratio of Physical Quantities}
\label{subsubsec:invariant_ratio}

Regardless of the chosen system of units, $\alpha$ remains constant as a \textbf{ratio} of fundamental lengths:

\begin{equation}
	\alpha_{\text{EM}} = \frac{r_e}{\lambda_C/2\pi} = \frac{\text{Classical electron radius}}{\text{Reduced Compton wavelength}}
	\label{eq:alpha_ratio_re}
\end{equation}

Similarly, the inverse ratio:

\begin{equation}
	\alpha_{\text{EM}}^{-1} = \frac{a_0}{\lambda_C/2\pi} = \frac{\text{Bohr radius}}{\text{Reduced Compton wavelength}} = 137.036...
	\label{eq:alpha_ratio_bohr}
\end{equation}

These ratios are \textbf{system-of-units invariant} -- they have the same numerical value in any consistent system of units, as the units cancel out in the ratio.

\subsubsection{As a System-Dependent Numerical Value}
\label{subsubsec:system_dependent}

Simultaneously, the numerical value of $\alpha$ depends on the choice of fundamental units:

\begin{itemize}
	\item \textbf{SI system}: $\alpha = \frac{e^2}{4\pi\varepsilon_0\hbar c} \approx 1/137$
	\item \textbf{Natural units}: $\alpha = 1$ (by suitable choice)
	\item \textbf{Gaussian units}: $\alpha = \frac{e^2}{\hbar c} \approx 1/137$
\end{itemize}

\subsection{The System Dependency of $\alpha$}
\label{subsec:alpha_system_dependency}

The numerical value $\alpha_{\text{EM}} = 1/137$ holds \textbf{in all systems that do not redefine $\alpha$ by convention} (SI, Gaussian, Heaviside-Lorentz); only the T0 setting turns it into $1$:

\begin{align}
	\text{SI system:} \quad &\alpha_{\text{EM}}^{\text{SI}} = \frac{e^2}{4\pi\varepsilon_0\hbar c} \approx \frac{1}{137.036} \\
	\text{Natural system of units:} \quad &\alpha_{\text{EM}}^{\text{nat}} = 1 \text{ (by suitable choice of units)}
\end{align}

\textbf{Transformation factor:}
\begin{equation}
	\frac{\alpha_{\text{EM}}^{\text{nat}}}{\alpha_{\text{EM}}^{\text{SI}}} = 137.036
\end{equation}

\subsection{The Natural System of Units with $\alpha = 1$}
\label{subsec:natural_units}

In a natural system of units that respects electromagnetic duality, we obtain:

\begin{itemize}
	\item $\hbar_{\text{nat}} = 1$ (quantum mechanical scale)
	\item $c_{\text{nat}} = 1$ (relativistic scale)
	\item $\varepsilon_{0,\text{nat}} = 1$ (electric constant)
	\item $\mu_{0,\text{nat}} = 1$ (magnetic constant)
	\item $e_{\text{nat}}^2 = 4\pi$ (elementary charge)
\end{itemize}

With these values, $\alpha = 1$ is verified in both the standard form and the dual form:

\begin{equation}
	\alpha = \frac{4\pi}{4\pi \cdot 1 \cdot 1 \cdot 1} = 1
\end{equation}

\subsection{The Resolution of the ``Mystery''}
\label{subsec:mystery_resolution}

The apparent mystification of $1/137$ arises from:

\begin{enumerate}
	\item \textbf{Confusion of two aspects}: The invariance of the ratios is conflated with the system-dependency of the numerical representation.
	
	\item \textbf{Treatment of the SI system as absolute}: The historically evolved SI units (meter, second, kilogram, ampere) force electromagnetic constants to take ``unnatural'' values.
	
	\item \textbf{Forgetting the construction of unit systems}: All unit systems are human constructs. Nature knows no preferred units.
	
	\item \textbf{Search for deeper meaning in conversion factors}: The number 137 has no deeper cosmic significance than, say, the factor 1609.344 between miles and meters.
\end{enumerate}

\subsection{The Anthropic Fallacy}
\label{subsec:anthropic_fallacy}

Typical anthropic arguments claim:
\begin{itemize}
	\item ``If $\alpha_{\text{EM}} = 1/200$ $\rightarrow$ no atoms $\rightarrow$ no life''
	\item ``If $\alpha_{\text{EM}} = 1/80$ $\rightarrow$ no stars $\rightarrow$ no life''
	\item ``Therefore, $\alpha_{\text{EM}} = 1/137$ is `fine-tuned' for life''
\end{itemize}

\textbf{The problem}: These arguments presuppose the SI system as absolute!

\textbf{In natural units}: $\alpha_{\text{EM}} = 1$ is perfectly natural and requires no fine-tuning whatsoever. The electromagnetic interaction has unit strength in the natural system of units, which respects the fundamental structure of quantum mechanics and relativity.

\subsection{Sommerfeld's Harmonic Imprinting}
\label{subsec:sommerfeld_harmonic}

An often overlooked historical aspect: In 1916, Arnold Sommerfeld actively searched for \textbf{harmonic ratios} in atomic spectra, guided by the philosophical conviction that nature follows musical principles.

His methodological approach:
\begin{enumerate}
	\item \textbf{Expectation} of musical ratios in quantum transitions
	\item \textbf{Calibration} of measurement systems to produce harmonic values
	\item \textbf{Definition} of $\alpha_{\text{EM}}$ based on harmonic spectroscopic adjustments
	\item \textbf{Attribution} of the resulting ratio to fundamental physics
\end{enumerate}

The apparent ``harmony'' in $\alpha_{\text{EM}}^{-1} = 137 \approx (6/5)^{27}$ is therefore not a cosmic discovery, but the result of Sommerfeld's harmonic expectations embedded into the definition of the unit system.

\subsection{Physical Interpretation}
\label{subsec:physical_interpretation}

In natural units, $\alpha = 1$ represents the perfect balance between:

\begin{itemize}
	\item \textbf{Electric field coupling} (via $\varepsilon_0$ with $c^{-1}$)
	\item \textbf{Magnetic field coupling} (via $\mu_0$ with $c^{+1}$)
	\item \textbf{Quantum mechanical scale} (via $\hbar$)
	\item \textbf{Relativistic scale} (via $c$)
\end{itemize}

The electromagnetic duality $\frac{1}{\varepsilon_0 c} = \mu_0 c$ ensures this perfect balance.

\subsection{Historical Warning: The Eddington Saga}
\label{subsec:eddington_warning}

Arthur Eddington (1882-1944) attempted to ``prove'' $\alpha_{\text{EM}} = 1/137$ from first principles and developed elaborate numerological theories. The result was entirely speculative and wrong -- a warning against mystifying system-dependent numbers.

However, modern analysis shows that the fine-structure constant is indeed derivable from fundamental electromagnetic vacuum constants and that $\alpha_{\text{EM}} = 1$ in natural units is not only possible but reveals the arbitrary nature of our choice of unit system.
	\section{The $\beta_{\text{T}}$ Parameter}
	\label{sec:beta_t}
	
	\subsection{Theoretical Foundation from Field Theory}
	\label{subsec:beta_field_theory}
	
	In T0-natural units $\betaTnat = 1$. The T0 model connects this value with the field-theoretic relation \cite{pascher_derivation_beta_2025}:
	
	\begin{equation}
		\beta_{\text{T}} = \frac{\lambdah^2 v^2}{16\pichar^3 m_h^2 \xipar} \approx 1
		\label{eq:beta_t_field_theory}
	\end{equation}
	
	Since the Higgs expression matches $\xipar$ only as a consistency check (deviation of about $2.3\,\%$, Doc.~385), this relation supports $\beta_{\text{T}} = 1$ in natural units but does not fix the value exactly.
	
	\subsection{No Measured SI Value}
	\label{subsec:beta_circularity}
	
	A determination of $\betaTSI$ via a wavelength-dependent redshift $z(\lambda) = z_0\left(1 + \beta_{\text{T}} \ln(\lambda/\lambda_0)\right)$ does not hold, because the redshift is achromatic in the T0 framework: the fractal path lengthening stretches all wavelengths by the same factor $1+z$, independent of $\lambda$; the scale of this factor is calibrated via the measured $H_0$ (P39). An observationally determined SI value of $\beta_{\text{T}}$ and a conversion factor between the systems therefore do not arise; unlike $\alpha_{\text{EM}}$, $\beta_{\text{T}}$ is no example of a measured, system-dependent numerical value.
	
	\section{The Wien Constant $\alpha_{\text{W}}$}
	\label{sec:alpha_w}
	
	\subsection{Mathematical vs. Conventional Values}
	\label{subsec:wien_values}
	
	Wien's displacement law gives:
	
	\begin{align}
		\text{SI system:} \quad &\alphaWSI = 2.8977719...\times10^{-3}~\text{m\,K} \\
		\text{T0 system:} \quad &\alphaWnat = 2\pi/x_W = 1.2655 \quad (x_W = 4.965114)
	\end{align}
	
	\textbf{Conversion:}
	$$\alphaWSI = \alphaWnat\cdot\frac{\hbar c}{k_B} = 1.2655 \times 2.2899\times10^{-3}~\text{m\,K}$$
	In $\hbar = c = k_B = 1$ Wien's constant is fixed as $2\pi/x_W = 1.2655$ and cannot additionally be set to 1. There is therefore no dimensionless conversion factor between the systems for $\alpha_{\text{W}}$.
	
	\section{Parameter Comparison Table}
	\label{sec:parameter_comparison}
	
	\begin{table}[htbp]
		\centering

\adjustbox{max width=\textwidth}{%
\begin{tabular}{lcccc}
			\toprule
			\textbf{Parameter} & \textbf{SI Value} & \textbf{T0-nat Value} & \textbf{Ratio} & \textbf{Factor} \\
			\midrule
			$\xipar$ (electron) & $8.4 \times 10^{-23}$ & $8.4 \times 10^{-23}$ & 1 & $10^{0}$ \\
			$\alpha_{\text{EM}}$ & $7.3 \times 10^{-3}$ & $1$ & 137 & $10^{2.1}$ \\
			$\beta_{\text{T}}$ & -- (no measured value) & $1$ & -- & -- \\
			$\alpha_{\text{W}}$ & $2.898\times10^{-3}$~m\,K & $1.2655$ & $\hbar c/k_B$ & -- \\
			\bottomrule
		\end{tabular}%
}
		\caption{Systematic parameter differences between unit systems}
		\label{tab:parameter_comparison}
	\end{table}
	
	\textbf{$\alpha_{\text{EM}}$ shows about 2 orders of magnitude difference between systems; $\beta_{\text{T}}$ has no measured SI value (achromatic redshift); for $\alpha_{\text{W}}$ there is no dimensionless conversion factor (conversion via $\hbar c/k_B$); $\xipar$ is the same in both.}
	
	\section{Yukawa Parameters: Variable and System-Dependent}
	\label{sec:yukawa_parameters}
	
	\subsection{The Hierarchy of Yukawa Couplings}
	\label{subsec:yukawa_hierarchy}
	
	In the Standard Model, Yukawa couplings vary dramatically:
	
	\begin{table}[htbp]
		\centering
		\adjustbox{max width=\textwidth}{%
\begin{tabular}{lc}
			\toprule
			\textbf{Particle} & \textbf{$y_i$ (SI system)} \\
			\midrule
			Electron & $2.94 \times 10^{-6}$ \\
			Muon & $6.09 \times 10^{-4}$ \\
			Tau & $1.03 \times 10^{-2}$ \\
			Up quark & $1.27 \times 10^{-5}$ \\
			Top quark & $1.00$ \\
			Bottom quark & $2.25 \times 10^{-2}$ \\
			\bottomrule
		\end{tabular}%
}
		\caption{Yukawa coupling hierarchy (5 orders of magnitude variation)}
		\label{tab:yukawa_hierarchy}
	\end{table}
	
	\subsection{Transformation Uncertainty}
	\label{subsec:yukawa_transformation}
	
	The transformation of Yukawa parameters between systems requires careful consideration of the Higgs mechanism. The general form would be:
	
	$$y_{i,\text{nat}} = y_{i,\text{SI}} \times T_{\text{Yukawa}}$$
	
	where $T_{\text{Yukawa}}$ depends on the transformation of Higgs vacuum expectation value and particle masses.
	
	\subsection{Consistency Requirements}
	\label{subsec:yukawa_consistency}
	
	The Higgs mechanism requires:
	$$m_h^2 = 2\lambdah v^2$$
	
	For transformation consistency:
	$$T_m^2 = T_\lambda \times T_v^2$$
	
	This gives:
	$$y_{i,\text{nat}} = y_{i,\text{SI}} \times \sqrt{T_\lambda}$$
	
	\textbf{However, $T_\lambda$ requires detailed specification of the T0-natural unit system transformation rules.}
	
	\section{Universal Warning: No Direct Parameter Transfer}
	\label{sec:universal_warning}
	
	\subsection{The Systematic Problem}
	\label{subsec:systematic_problem}
	
	\begin{warning}
		\textbf{EVERY parameter symbol in T0-model documents may have different values than in SI system calculations!}
	\end{warning}
	
	\textbf{Concrete danger zones:}
	
	\begin{align}
		G_{\text{nat}} &= 1 \quad \text{vs.} \quad G_{\text{SI}} = 6.674 \times 10^{-11} \text{ m}^3/(\text{kg} \cdot \text{s}^2) \\
		\alpha_{\text{EM,nat}} &= 1 \quad \text{vs.} \quad \alpha_{\text{EM,SI}} = 1/137 \\
		e_{\text{nat}} &= 2\sqrt{\pichar} \quad \text{vs.} \quad e_{\text{SI}} = 1.602 \times 10^{-19} \text{ C}
	\end{align}
	
	\textbf{Direct transfer leads to errors of factors $10^2$ to $10^{11}$!}
	
	\subsection{Required Transformation Protocol}
	\label{subsec:transformation_protocol}
	
	For every parameter, explicitly specify:
	
	\begin{enumerate}
		\item \textbf{Which unit system} is being used
		\item \textbf{How transformation occurs} between systems
		\item \textbf{Which factors must be considered}
		\item \textbf{Which consistency conditions} must be satisfied
	\end{enumerate}
	
	\textbf{Example of complete specification:}
	\begin{tcolorbox}[colback=red!5!white,colframe=red!75!black,title=Parameter Specification Template]
		\textbf{Parameter:} Fine structure constant $\alpha_{\text{EM}}$ \\
		\textbf{SI value:} $\alphaEMSI = 1/137.036$ \\
		\textbf{T0 value:} $\alphaEMnat = 1$ \\
		\textbf{Transformation:} $\alphaEMnat = \alphaEMSI \times 137.036$ \\
		\textbf{Consistency:} Dimensional analysis verified \\
		\textbf{Usage:} Specify system before calculation
	\end{tcolorbox}
	
	\subsection{Experimental Prediction Guidelines}
	\label{subsec:experimental_guidelines}
	
	\textbf{For QED calculations:}
	\begin{align}
		\text{WRONG:} \quad &\alpha_{\text{EM}} = 1 \text{ from T0-model directly in SI formulas} \\
		\text{CORRECT:} \quad &\alphaEMSI = 1/137 \text{ with transformation to } \alphaEMnat = 1
	\end{align}
	
	\textbf{For gravitational calculations:}
	\begin{align}
		\text{WRONG:} \quad &G = 1 \text{ from T0-model directly in Newton's formulas} \\
		\text{CORRECT:} \quad &G_{\text{SI}} = 6.674 \times 10^{-11} \text{ with transformation to } G_{\text{nat}} = 1
	\end{align}
	
	\section{The Circularity Resolution}
	\label{sec:circularity_resolution}
	
	\subsection{Apparent vs. Real Circularity}
	\label{subsec:apparent_real_circularity}
	
	The circularity problem that seemed to plague T0-model parameter determination is resolved by recognizing:
	
	\begin{enumerate}
		\item \textbf{No real circularity exists} within each consistent system
		\item \textbf{Both SI and T0 systems are internally consistent}
		\item \textbf{The apparent contradiction} arose from comparing parameters across different systems
		\item \textbf{Proper transformation} eliminates all apparent inconsistencies
	\end{enumerate}
	
	\subsection{System Consistency Verification}
	\label{subsec:system_consistency}
	
	\textbf{SI system consistency:}
	$$\Rzero = \frac{m_e c \left(\alphaEMSI\right)^2}{2h} \quad \checkmark \text{ (experimentally verified to 0.000001\%)}$$
	
	\textbf{T0 system consistency:}
	$$\hbar = c = G = k_B = \alpha_{\text{EM}} = \beta_{\text{T}} = 1 \quad \checkmark \text{ (by construction)}$$
	
	\textbf{Both systems work perfectly within their own frameworks!}
	
	\section{Implications for T0-Model Testing}
	\label{sec:testing_implications}
	
	\subsection{System-Specific Predictions}
	\label{subsec:system_specific_predictions}
	
	Experimental tests must clearly specify which parameter system is used:
	
	\begin{table}[htbp]
		\centering
		\adjustbox{max width=\textwidth}{%
\begin{tabular}{lcc}
			\toprule
			\textbf{Test Type} & \textbf{SI-based Prediction} & \textbf{T0-based Prediction} \\
			\midrule
			QED anomaly & $a_e \propto \alphaEMSI = 1/137$ & $a_e \propto \alphaEMnat = 1$ \\
			Galaxy rotation & $v^2 \propto \xipar \sim 10^{38}$ & $v^2 \propto \xipar \sim 10^{38}$ (same) \\
			\bottomrule
		\end{tabular}%
}
		\caption{System-specific experimental predictions}
		\label{tab:system_predictions}
	\end{table}
	
	\subsection{Transformation Validation}
	\label{subsec:transformation_validation}
	
	The transformation factors can be validated by checking:
	
	\begin{enumerate}
		\item \textbf{Dimensional consistency} in both systems
		\item \textbf{Known limits} are reproduced correctly
		\item \textbf{Ratios remain invariant} between systems
		\item \textbf{Internal consistency} of each system
	\end{enumerate}
	% Bibliographie
\bibliographystyle{natbib}
\begin{thebibliography}{99}
	\bibitem{pascher_derivation_beta_2025}
	J. Pascher.
	\newblock Derivation of Beta in the T0 Framework.
	\newblock T0-Theory Document Collection, 2025.
	\newblock URL: \url{https://github.com/jpascher/T0-Time-Mass-Duality/blob/main/2/pdf/093_DerivationVonBeta_En.pdf}
	(verifizierter Link existiert).
\end{thebibliography}